\documentclass[12pt]{article}
\usepackage[utf8]{inputenc}
\usepackage[T1]{fontenc}
\usepackage[margin=1in]{geometry}
\usepackage{setspace}
\usepackage{xcolor}
\usepackage[colorlinks=true, urlcolor=blue!60!black, linkcolor=blue!60!black]{hyperref}
\setlength{\parindent}{0em}
\setlength{\parskip}{0.45em}
\renewcommand{\baselinestretch}{1.0}
\begin{document}
\begin{center}
{\Large \textit{Youngstown Sheet \& Tube Co. v. Sawyer}, 343 U.S. 579 (1952)\par}
\vspace{0.3em}
{\normalsize Case Brief --- One-Page Edition\par}
\end{center}

\textbf{Memory Jogger.} The ``Steel Seizure Case'' and Justice Jackson's three zones of presidential power (\href{https://supreme.justia.com/cases/federal/us/343/579/}{343 U.S. 579}).

\textbf{Facts.} Facing a steel strike during the Korean War, Truman bypassed the Taft-Hartley Act and ordered Commerce Secretary Sawyer to seize the mills, citing no statute. Congress did nothing.

\textbf{Procedural History.} The district court enjoined the seizure; the Court granted certiorari within weeks.

\textbf{Votes.} Affirmed, 6--3: Black for the Court, with five separate concurrences; Vinson dissented for three.

\textbf{Holding.} Presidential power must come from Congress or the Constitution; no Article II clause permits seizing mills to settle a labor dispute, a lawmaking function Congress had withheld.

\textbf{Reasoning.} Black: the Take Care Clause confines the President to recommending and vetoing laws. Jackson: power is ``at its maximum'' with Congress, in a ``zone of twilight'' when Congress is silent, and ``at its lowest ebb'' against Congress's will, where Truman stood. Frankfurter required long acquiesced-in practice; Vinson mocked a ``messenger-boy concept of the Office.''

\textbf{Examples.} \emph{Same-Side:} seizing chip plants against a labor statute. \emph{Opposite-Side:} IEEPA asset freezes (\textit{Dames \& Moore}). \emph{Fence-Sitter:} an emergency grid-security order amid congressional silence.

\textbf{Critique.} Formalists praise Black; others see Vinson's ``completion power'' vindicated or question Zone 3's assumption of ``categorical congressional supremacy.''

\textbf{Key Quotation.} ``While the Constitution diffuses power the better to secure liberty, it also contemplates that practice will integrate the dispersed powers into a workable government. It enjoins upon its branches separateness but interdependence, autonomy but reciprocity. Presidential powers are not fixed but fluctuate depending upon their disjunction or conjunction with those of Congress.'' --- Justice Jackson, concurring, 343 U.S. at 635.

\end{document}
